\chapter{wpw 系统需求分析}

第 3 章提出了面向 AI 辅助开发的上下文工程方法，包括流程约束架构、能力-代码双层知识模型、图谱驱动的任务上下文生成和 Spec 增量知识演化。方法要发挥作用，需要一套可运行的系统把它承载下来：阶段约束要落在可检查的状态上，知识模型要落在可持久化的文件上，上下文生成要落在可调用的命令上。本章分析这个系统——wpw——的需求，说明系统的使用者和使用场景，明确系统应具备的功能、应满足的非功能要求以及建设边界，为第 5 章的设计与实现划定范围。


\section{系统概述与使用者}
wpw 是第 3 章方法的工程载体。它的定位不是替代大模型或智能体的运行框架，而是在智能体已经具备代码生成和工具调用能力的前提下，为长期演化项目提供可维护的过程制品、项目知识和任务级上下文。系统以命令行工具和智能体阶段命令的形式运行，过程制品、能力规范和知识图谱均以本地文件保存，不要求部署独立的数据库服务。

系统有两类使用者。第一类是开发者，负责提出需求、确认各阶段制品、在技术探索的候选方案中拍板、验证实施结果并决定是否归档。第二类是 AI 编程智能体，负责生成阶段制品、调用系统命令、按任务请求上下文并实施代码修改。两类使用者的职责需要明确分开：智能体的生成行为不能直接改写项目状态，项目状态的每一次变化都要经过确定性程序记录。这样，模型可以被替换，会话可以中断，而项目状态和已积累的知识仍然完整可查。

系统的典型使用场景是长期演化的 Web 项目。项目接入系统时，先对现有代码做一次全量解析，建立模块、文件和代码元素的结构索引，并结合已有需求规范形成初始知识图谱。此后每次需求变更作为一个独立的过程实例进入系统，经历业务需求、产品需求、技术设计、开发计划、测试和实施的完整流程；实施过程中，智能体按任务向图谱请求与当前需求相关的代码上下文；需求完成归档时，已确认的需求增量被沉淀到能力规范中，图谱随之增量更新。系统的价值不在单次任务，而在这一周期随项目演化反复执行后，知识仍然可用、上下文成本仍然可控。


\section{功能需求}
功能需求对应第 3 章方法的四个组成部分：流程约束架构要求系统提供工作流约束能力，双层知识模型要求系统提供图谱构建能力，任务上下文生成方法要求系统提供上下文生成能力，Spec 增量知识演化要求系统提供增量维护能力。表4-1 汇总了四个功能域下的主要需求项，后续小节分别说明。

\begin{table}[htbp]
\centering
\small
\caption{表4-1 wpw 功能需求汇总\\Table 4-1 Summary of Functional Requirements of wpw}
\begin{tabularx}{\textwidth}{|X|X|X|}
\hline
\textbf{功能域} & \textbf{需求项} & \textbf{需求描述} \\
\hline
工作流约束 & 需求实例管理 & 每次需求变更创建独立工作区，校验名称格式与全局重名，初始化各阶段状态 \\
工作流约束 & 阶段推进控制 & 按阶段前置关系实施门禁：强依赖未完成不得进入下一阶段，可选阶段可显式跳过并留痕 \\
工作流约束 & 阶段制品支持 & 提供各阶段模板，检查前序制品是否齐备，制品确认后记录状态 \\
工作流约束 & 任务实施支持 & 以开发计划中的任务清单为实施来源，支持任务级进度标记和中断后按显式状态恢复 \\
工作流约束 & 归档处理 & 需求完成后迁移工作区至归档目录，并衔接能力沉淀与图谱更新 \\
图谱构建 & 源码结构解析 & 解析多语言源代码，建立模块、文件、代码元素三层结构及导入、调用、继承关系 \\
图谱构建 & 能力规范解析 & 从能力规范文件中提取业务能力名称、用途、需求和场景，形成能力层节点 \\
图谱构建 & 业务—代码映射 & 汇聚文档提取、命名匹配、语义相似度和 Git 历史等证据，生成能力与代码之间的映射边 \\
图谱构建 & 统一存储 & 以统一数据模型保存节点和边，提供持久化文件、索引和元数据 \\
上下文生成 & 任务锚点选择 & 支持自然语言任务的语义锚点检索，也支持直接指定模块、文件或元素锚点 \\
上下文生成 & 子图扩展与裁剪 & 从锚点沿图边扩展并裁剪，受深度、边权重和节点数约束 \\
上下文生成 & 序列化输出 & 按层级组织子图内容，以紧凑形式输出结构化上下文 \\
上下文生成 & 预算控制 & 以 Token 预算为输入约束，超预算时逐级降低信息粒度，无法满足时明确报告 \\
增量维护 & 变更识别 & 比对文件哈希快照，识别新增、修改和删除的源文件与能力规范 \\
增量维护 & 局部更新 & 只重建受影响的节点和边，未变化部分保留；版本不兼容时自动转为全量构建 \\
增量维护 & 能力沉淀 & 归档时将已确认的需求增量合并到既有能力规范，或按统一结构新建能力规范 \\
\hline
\end{tabularx}
\end{table}


\subsection{工作流约束需求}
工作流约束是流程约束架构对系统的要求。系统需要把一次需求变更当作一个独立的过程实例来管理，而不是依赖一段连续对话：实例创建时即固定需求名称、工作目录和初始状态，后续命令凭实例标识定位，而不需要从聊天记录推断当前处于哪个阶段。

阶段推进需要可检查的门禁。业务需求、产品需求、技术设计、开发计划和实施构成主链，技术探索和测试作为可选环节。进入任一阶段前，系统先检查其强依赖阶段的制品是否已确认；未满足时阻止进入并说明缺失项。可选阶段允许跳过，但跳过必须是显式记录的决定；技术探索已形成候选方案而未拍板时，系统应在后续阶段给出可见警告，而不是静默放行。

实施阶段以开发计划中的任务清单为单位推进。系统需要支持任务级进度标记，使已完成和剩余任务可随时查询；会话中断后，智能体应能凭状态文件和阶段制品恢复现场，继续执行剩余任务。需求完成并经开发者验证后，系统提供归档处理，将过程工作区迁移到归档目录，并把这次变更中确认的需求增量交给知识维护流程。


\subsection{图谱构建需求}
图谱构建是双层知识模型对系统的要求。系统需要解析项目源代码，建立模块、文件、代码元素三个结构层级，并识别导入、调用和继承等关系。项目接入时不应假设单一语言或单一框架，常见的 Web 项目形态，如前后端分离、状态管理单元和单文件组件，都应在解析范围内；构建前还需按文件类型和内容特征预检，避免把构建产物、依赖目录和文档资料混入代码结构层。

能力层来自能力规范文件。系统需要从规范中解析出业务能力的名称、用途、需求和场景，形成能力节点，并汇聚多源证据生成能力与代码结构之间的映射边：文档中提取的显式关联、命名上的对应、语义上的相似以及版本历史中能力与代码的共同变更，都可以作为证据；证据经合成后形成带权重的映射边，并记录主要证据来源以备检查。某一证据来源不可用时，其余证据仍应能形成映射，而不是导致构建失败。

图谱的持久化需要统一的数据模型。不同层级、不同语言解析器产生的节点和边应以同一格式保存、索引和传递，使上下文生成模块无需关心节点的来源语言。图谱文件、索引和元数据都应落在本地文件中，并记录构建时间、节点与边的规模和配置摘要，便于核查与重建。


\subsection{上下文生成需求}
上下文生成是任务级子图生成方法对系统的要求。系统接受两类任务输入：智能体以自然语言描述任务时，系统通过语义检索找到候选锚点；调用方已知相关模块、文件或元素标识时，可以直接指定锚点，跳过语义检索。

从锚点出发，系统沿图谱的结构边和业务映射边扩展，得到与任务相关的候选子图。扩展必须受深度、边权重和节点数量的约束，并按语义相关性、结构重要性和到锚点的距离裁剪，避免把无关代码带入上下文。裁剪后的子图按能力、模块、文件、元素的层次组织输出，锚点保留较完整的属性与签名，较远节点只保留名称、类型和必要关系，以紧凑形式减少重复 Token。

Token 预算是上下文生成的输入约束，而不是隐藏在提示词中的经验参数。调用方给出预算后，系统应在超预算时逐级降低信息粒度，并在各级降级中始终保留锚点；如果连只含锚点的最小上下文仍超出预算，系统应明确报告预算不足，而不是静默截断，使调用方能够察觉并调整任务粒度。


\subsection{增量维护需求}
增量维护是 Spec 增量知识演化对系统的要求。长期项目中，每次代码修改都全量重建图谱会让知识维护本身成为负担。系统需要保存源码文件的哈希快照，在更新时识别新增、修改和删除的文件，只重建受影响文件的节点和关系，未变化的部分保留。能力规范发生变化时，系统重新解析并与既有能力节点比对，处理能力的新增、修改和删除，并重新计算受影响的映射边。

更新以数据结构版本兼容为前提。尚未存在历史图谱，或发现主版本不兼容时，系统应自动转为全量构建，并提供显式的重建命令以应对结构性重构。能力沉淀与归档流程衔接：需求归档时，系统依据已确认的需求制品和实施结果，将需求增量合并到既有能力规范，或按统一结构创建新的能力规范，随后触发图谱更新，使本次变更的知识进入项目知识库。


\section{非功能需求}
功能需求规定系统做什么，非功能需求规定系统以什么方式存在和运行。wpw 服务于长期演化项目，其非功能需求围绕部署成本、过程可靠性和行为可验证性展开。表4-2 汇总了主要的非功能需求。

\begin{table}[htbp]
\centering
\small
\caption{表4-2 wpw 非功能需求汇总\\Table 4-2 Summary of Non-Functional Requirements of wpw}
\begin{tabularx}{\textwidth}{|X|X|}
\hline
\textbf{类别} & \textbf{需求描述} \\
\hline
部署与存储 & 零数据库：图谱、索引和状态均以本地文件保存，随项目目录迁移，不要求独立服务 \\
可靠性 & 状态文件与阶段文档构成可恢复的过程快照，会话中断后从显式状态恢复；文件写入采用临时文件加原子重命名，避免半成品文件 \\
可降级性 & 图谱未构建、向量索引缺失或本地语义模型不可用时，回退到显式锚点或人工读取文件，不构成流程硬阻塞 \\
可扩展性 & 阶段规则集中于统一的模式定义，解析器按语言扩展，阶段模板允许项目级覆盖 \\
可验证性 & 命令输出结构化结果，测试以持久化状态和文件为判据，不以单次模型回答为判据 \\
环境约束 & 运行环境为 Node.js 18 及以上；版本信息、配置摘要和图谱模式版本可记录、可锁定 \\
\hline
\end{tabularx}
\end{table}

部署上，系统不引入独立数据库。图谱以行式文本文件保存，向量索引以二进制文件保存，连同状态文件都放在项目目录内。这样项目迁移、备份和版本管理都不需要额外的数据导出步骤，也使实验环境易于复现。

可靠性上，系统的状态变更必须由确定性程序完成并留下记录。需求名称、各阶段状态、任务进度和归档位置共同构成过程快照；智能体与开发者的会话可以中断、重来，项目状态不因此丢失。图谱和元数据文件的写入应先写临时文件再重命名，避免写入中断留下不完整文件。

可降级性上，知识服务是流程的增强而非前提。图谱不存在、向量索引缺失或锚点未命中时，系统给出提示并回退到显式锚点查询或人工读取文件，开发流程可以继续推进。这一要求直接影响设计：所有依赖图谱的命令都需要定义无图谱时的行为。

可验证性上，系统命令应输出结构化结果，功能测试以命令输出、状态转换和持久化文件为判据。这样系统行为可以被脚本化检查，也为第 6 章的系统测试和实验记录提供统一的数据来源。


\section{系统约束与建设边界}
系统建设受三方面约束。技术环境方面，系统以 Node.js 和 TypeScript 实现，源码解析以 tree-sitter 为基础，当前适配 TypeScript、JavaScript、Vue 单文件组件和 Java；未适配的语言不进入图谱，但不影响工作流的使用。项目形态方面，系统面向有阶段制品和版本管理的软件项目，长期演化、多轮需求迭代的项目更能体现知识积累的价值；一次性、短周期项目也可以使用工作流，但图谱构建的投入相对收益有限。运行方式方面，系统以命令行工具和智能体阶段命令的形式挂载在智能体运行框架之上，不要求改造智能体本身。

建设边界同时是本文的研究边界在系统层面的落实。其一，系统不改动基础模型，不进行训练或微调；图谱构建中的多源证据只是工程化的关联依据，证据的强弱不影响模型本身的能力。其二，系统不替代智能体的运行框架，智能体的推理与执行仍由其运行框架负责，系统只管理过程制品、项目知识和任务上下文。其三，系统不评价下游任务的算法效果；后续实验载体中的推荐算法仅用于构成真实的开发任务，其准确率不是研究变量。其四，系统输出的是与任务相关的结构化上下文，帮助智能体和开发者更快定位相关代码，但不替代开发者对关键源码的精读与判断；能力沉淀的质量也依赖开发者对需求制品和合并结果的审阅。


\section{本章小结}
本章分析了 wpw 系统的需求。系统面向开发者和 AI 编程智能体两类使用者，服务于长期演化项目的完整开发周期。功能需求覆盖工作流约束、图谱构建、上下文生成和增量维护四个方面，分别承载第 3 章方法的四个组成部分；非功能需求要求系统零数据库部署、状态可恢复、服务可降级、行为可验证；系统约束与建设边界明确了系统的技术环境、适用项目形态，以及不改动基础模型、不替代智能体运行框架、不评价下游算法效果的边界。第 5 章将按照这些需求，说明系统的三层架构以及工作流、图谱构建、上下文生成和增量维护四个模块的设计与实现。
